\section{Proposed Method}
\label{sec:method}
\subsection{Overview}
Describe the complete procedure from input to output. Refer to
Fig.~\ref{fig:overview} and explain the role of each component.
Separate your own contribution from standard components or reused methods.

\subsection{Main Component}
Explain the central mechanism in enough detail for another student to
understand and implement it. Define any parameters, dependencies, and
design choices. Explain why those choices follow from the problem.

If an optimization objective is relevant, use the following layout as an
example and replace the expression with your actual objective:
\begin{equation}
  \theta^{\star}=\arg\min_{\theta}
  \frac{1}{N}\sum_{i=1}^{N}\ell\bigl(f_{\theta}(x_i),y_i\bigr).
  \label{eq:objective}
\end{equation}
Here, $N$ is the number of examples, $(x_i,y_i)$ is an input--target pair,
and $\ell$ is a loss function. This generic equation is not a proposed
research contribution and may be deleted.

\subsection{Implementation Details}
Report the settings needed to reproduce your method, such as preprocessing,
parameter selection, software dependencies, and computational resources.
For a conceptual proposal, distinguish planned steps from implemented ones.

\subsection{Component Design and Information Flow}
Figure~\ref{fig:component} provides a closer view of a single component.
Use this figure to explain an internal mechanism that would be too detailed
for the system overview. Name the information entering each stage, describe
the operation applied to it, and specify what is passed to the next stage.
Keep terminology consistent between the diagram, equations, and text.

\begin{figure}[!t]
  \centering
  \fbox{\begin{minipage}[c][43mm][c]{\dimexpr\columnwidth-2\fboxsep-2\fboxrule\relax}
    \centering
    \textbf{DUMMY: COMPONENT DETAIL}\\[4mm]
    Input representation\\[2mm]
    $\downarrow$\\[2mm]
    \fbox{Stage A: transform}\quad$\rightarrow$\quad\fbox{Stage B: refine}\\[2mm]
    $\downarrow$\\[2mm]
    Output representation
  \end{minipage}}
  \caption{Mechanism diagram placeholder. Show how one component works
  internally, including intermediate representations and operations.
  This figure complements the full-system overview in Fig.~\ref{fig:overview}.}
  \label{fig:component}
\end{figure}


Explain how intermediate representations are constructed. Describe whether
they are computed independently for each input or depend on earlier inputs.
If the method maintains a state, define its initialization and update rule.
If it combines multiple sources of information, explain how their relative
influence is determined. State what happens when an expected input is
unavailable, incomplete, or outside the assumed operating range.

Describe the interface between your proposed component and the rest of the
system. A reader should be able to distinguish the component's inputs and
outputs from implementation conveniences. Identify any quantities that are
available during development but unavailable at evaluation time. Explain
how the evaluation prevents information from the test set from influencing
the method or its settings.

\subsection{Design Rationale and Alternatives}
For each major design decision, describe the alternative that a reasonable
reader might consider. Explain why the selected option is appropriate for
the stated problem. These explanations should be testable where possible:
identify which experimental comparison could support or weaken each claim.
Avoid presenting a convenient implementation choice as a demonstrated
scientific advantage.

\subsection{Reproducibility and Computational Cost}
Describe the software environment, input preparation, initialization,
termination criteria, and sources of randomness. Report how settings are
selected and which settings remain fixed across experiments. Distinguish
one-time preparation costs from per-input costs, and specify whether
reported timing includes loading data and producing the final output.
